\documentclass[aspectratio=169,11pt]{beamer}

\usetheme{Madrid}
\usecolortheme{default}
\usefonttheme{professionalfonts}
\setbeamertemplate{navigation symbols}{}
\setbeamertemplate{footline}[frame number]

\usepackage{amsmath, amssymb, booktabs, graphicx}

\newcommand{\E}{\mathbb{E}}
\newcommand{\Prb}{\mathbb{P}}

\title{Potential Outcomes, Identification Logic, and Selection-On-Observables}
\subtitle{Empirical Methods --- Week 1}
\author{Teaching deck draft}
\date{}

\begin{document}

\begin{frame}
  \titlepage
\end{frame}

\begin{frame}{Week 1 question}
\begin{itemize}
    \item What does it mean to identify a causal effect?
    \item What is observed, what is missing, and what assumption links them?
    \item When can rich observables support a credible comparison?
\end{itemize}

\medskip
The lecture uses LaLonde, Dehejia--Wahba, Rosenbaum--Rubin, and sensitivity-style work as the paper spine.
\end{frame}

\begin{frame}{Potential outcomes}
\[
Y_i = D_iY_i(1) + (1-D_i)Y_i(0)
\]

\[
\tau_i = Y_i(1)-Y_i(0)
\]

\begin{itemize}
    \item Treated units reveal $Y_i(1)$ and hide $Y_i(0)$.
    \item Untreated units reveal $Y_i(0)$ and hide $Y_i(1)$.
    \item Identification is an argument for replacing the missing potential outcome with a credible comparison.
\end{itemize}
\end{frame}

\begin{frame}{Target parameters}
\[
ATE = \E[Y(1)-Y(0)]
\]
\[
ATT = \E[Y(1)-Y(0)\mid D=1]
\]
\[
ATC = \E[Y(1)-Y(0)\mid D=0]
\]

\begin{itemize}
    \item ATE: average effect for the population represented by the sample.
    \item ATT: effect for units who actually received treatment.
    \item ATC: effect for units who are currently untreated.
    \item Support restrictions and weights can quietly change the parameter.
\end{itemize}
\end{frame}

\begin{frame}{Selection bias decomposition}
\[
\E[Y\mid D=1] - \E[Y\mid D=0]
=
\underbrace{\E[Y(1)-Y(0)\mid D=1]}_{ATT}
+
\underbrace{\E[Y(0)\mid D=1]-\E[Y(0)\mid D=0]}_{\text{selection bias}}
\]

\begin{itemize}
    \item The raw comparison is causal only if untreated potential outcomes are comparable.
    \item LaLonde shows how badly this can fail with the wrong nonexperimental comparison group.
    \item The comparison group is the design.
\end{itemize}
\end{frame}

\begin{frame}{Conditional independence and overlap}
\[
\{Y_i(1),Y_i(0)\} \perp D_i \mid X_i
\]

\[
0 < \Prb(D_i=1\mid X_i=x) < 1
\quad \text{for relevant } x
\]

\begin{itemize}
    \item Conditional independence: observables capture the assignment channels that matter for potential outcomes.
    \item Overlap: treated and untreated units coexist in the relevant covariate region.
    \item The propensity score balances observables under the identifying assumption; it does not create the assumption.
\end{itemize}
\end{frame}

\begin{frame}{Selection-on-observables is strongest when}
\begin{itemize}
    \item treatment assignment is institutionally understood;
    \item pre-treatment outcomes and histories are observed;
    \item covariates are measured before treatment;
    \item untreated comparisons come from the same economic environment;
    \item common support is strong;
    \item the main selection margins are observed rather than private or latent.
\end{itemize}
\end{frame}

\begin{frame}{Regression, matching, weighting}
\small
\begin{tabular}{p{0.25\linewidth} p{0.34\linewidth} p{0.30\linewidth}}
\toprule
Approach & Design idea & Main risk \\
\midrule
Regression adjustment & Model $\E[Y\mid D,X]$ and compare adjusted outcomes & Functional-form extrapolation \\
Matching & Compare treated units to similar untreated units & Bad matches or support loss \\
Weighting & Reweight by $e(X)=\Pr(D=1\mid X)$ & Extreme weights when overlap is weak \\
\bottomrule
\end{tabular}

\medskip
All three operationalize the same identifying claim: conditional comparability on observed pre-treatment information.
\end{frame}

\begin{frame}{IPW and doubly robust logic}
\[
ATE =
\E\left[
\frac{DY}{e(X)} -
\frac{(1-D)Y}{1-e(X)}
\right]
\]

\[
\E\left[
\hat m_1(X)-\hat m_0(X)
+
\frac{D(Y-\hat m_1(X))}{\hat e(X)}
-
\frac{(1-D)(Y-\hat m_0(X))}{1-\hat e(X)}
\right]
\]

\begin{itemize}
    \item Weighting uses the treatment model to reconstruct a target population.
    \item Doubly robust estimators combine treatment and outcome models.
    \item Neither relaxes the need for a credible observables-based assignment story.
\end{itemize}
\end{frame}

\begin{frame}{Papers as design lessons}
\small
\begin{tabular}{p{0.28\linewidth} p{0.58\linewidth}}
\toprule
Paper role & Design lesson \\
\midrule
LaLonde & Wrong comparison groups can miss the experimental benchmark badly. \\
Dehejia--Wahba & Matching can be credible when covariates, sample, and support are chosen deliberately. \\
Rosenbaum--Rubin & Propensity scores are balancing scores under conditional independence. \\
Altonji--Elder--Taber / Oster & Sensitivity work asks how strong omitted selection would need to be. \\
\bottomrule
\end{tabular}
\end{frame}

\begin{frame}{Diagnostics and robustness}
\begin{itemize}
    \item Balance: standardized differences before and after adjustment.
    \item Overlap: propensity-score support, trimming, and tails.
    \item Matching quality: exact variables, calipers, and substantively plausible pairs.
    \item Weighting stability: large weights and influential observations.
    \item Specification discipline: pre-treatment covariates and transparent model families.
    \item Sensitivity: omitted-variable robustness, placebo outcomes, negative controls.
\end{itemize}
\end{frame}

\begin{frame}{Research Lab design}
\begin{columns}[T,onlytextwidth]
\begin{column}{0.32\linewidth}
\textbf{Reproduce}
\begin{itemize}
    \item Synthetic LaLonde-style training data
    \item Naive, regression, matching, IPW, doubly robust estimates
\end{itemize}
\end{column}
\begin{column}{0.32\linewidth}
\textbf{Diagnose}
\begin{itemize}
    \item Balance
    \item Overlap
    \item Weight stability
    \item Estimand drift
\end{itemize}
\end{column}
\begin{column}{0.32\linewidth}
\textbf{Transfer}
\begin{itemize}
    \item Workplace training take-up
    \item Target parameter
    \item Selection story
    \item Omitted-variable threat
\end{itemize}
\end{column}
\end{columns}
\end{frame}

\begin{frame}{Takeaways}
\begin{itemize}
    \item Potential outcomes force the missing counterfactual into view.
    \item Selection bias is an economic assignment problem, not a regression nuisance.
    \item Selection-on-observables is credible only when observables capture the relevant selection story and support is real.
    \item Diagnostics and sensitivity analysis are part of the design, not an appendix ritual.
    \item Next week: experiments change the assignment mechanism, but not the need for estimand discipline.
\end{itemize}
\end{frame}

\end{document}
